\chapter{Introduction}
\label{ch:introduction}

High-resolution gridded precipitation data at 1~km resolution and below are
a critical input for hydrological modelling, flood and drought prediction,
climate impact assessment, and agricultural planning.
Producing these grids from observations is difficult: meteorological station
networks are sparse and spatially uneven, leaving large portions of any study
domain without direct measurements.
Converting discrete point observations into the continuous spatial fields
that downstream models require is the central problem of spatial precipitation
interpolation.

Precipitation's statistical properties make standard interpolation methods
inappropriate.
Unlike temperature or pressure, which are continuous and approximately
Gaussian, daily precipitation has a \textbf{mixed discrete-continuous} distribution.
A large fraction of all station-days record exactly zero (dry days), producing
a point mass at zero that cannot be represented by a single linear interpolator over a continuous field.
In the Central European study domain used in this thesis, dry days predominate
across the study period (overall dry-day fraction 51.9\%), though the proportion
varies by season and location, making the zero-inflation problem unavoidable.
Among wet days, precipitation amounts follow a heavy-tailed, right-skewed
distribution, with rare but intense events contributing disproportionately to
the total.

Both properties rule out naive application of standard kriging — a geostatistical method that estimates values at unobserved locations as a weighted average of neighbouring stations, assuming a stationary random field.

The core scientific problem addressed by this thesis is the transformation of
discrete, point-level daily precipitation measurements into continuous,
statistically coherent spatial fields at 1~km resolution over northern and
central Germany
(longitude $9.4^\circ$--$15.9^\circ$\,E, latitude $50.0^\circ$--$53.8^\circ$\,N,
EPSG:3035).
The study domain is served by the ReKIS regional climate network, comprising
2{,}458 stations with daily records spanning 1961--2023.

Three properties of the daily precipitation distribution require explicit treatment.
First, the \textbf{zero} inflation must be handled through either a two-stage model
that separates occurrence from amounts, or a unified framework that accommodates
the mixed distribution directly.
Second, the \textbf{non-stationarity of the mean}, caused by both the seasonal
cycle and orographic forcing, violates the stationarity assumption of standard
kriging (see Chapter~\ref{ch:theory}) and must be removed before spatial correlation can be modelled.
Third, the \textbf{right-skewness} of wet-day amounts requires a
variance-stabilising transformation if the kriging uncertainty estimates are
to be physically meaningful.

This thesis has three objectives:

\begin{enumerate}
    \item \textbf{Implement and validate Ordinary Kriging as a baseline.}
    A complete preprocessing and interpolation pipeline is developed to handle
    the statistical properties of daily precipitation that standard kriging
    cannot address directly.
    The pipeline is evaluated by cross-validation using a proper scoring
    rule that assesses both point accuracy and uncertainty calibration.

    \item \textbf{Assess the role of terrain elevation as an auxiliary predictor.}
    Elevation data from the Copernicus Digital Elevation Model are incorporated
    into the interpolation of climatological monthly normals.
    Cross-validation quantifies whether the orographic precipitation gradient
    in the study domain can be captured and whether doing so reduces
    interpolation error.

    \item \textbf{Contextualise results relative to machine learning alternatives.}
    Gradient boosting and Bayesian neural field approaches are implemented
    and compared as non-geostatistical alternatives.
    The geostatistical approach is the interpretable, physically motivated
    baseline; gradient boosting and the neural field model test whether
    data-driven flexibility improves on it.
\end{enumerate}

\textbf{Chapter~2} reviews the theoretical background of the three methods examined in this thesis.

\textbf{Chapter~3} describes the study domain and its data sources.

\textbf{Chapter~4} presents the implementation of the Ordinary Kriging pipeline.

\textbf{Chapter~5} presents the gradient boosting implementation for
spatial precipitation interpolation, including feature construction and
hyperparameter selection.

\textbf{Chapter~6} presents the Bayesian Neural Field implementation
applied to the same interpolation task.

\textbf{Chapter~7} compares the three methods on common evaluation metrics,
examines the gridded outputs, and discusses the trade-offs between
geostatistical interpretability and data-driven flexibility.

\textbf{Chapter~8} summarises the main findings and outlines directions for
future work.

\bigskip
\noindent\textbf{Code and data availability.}
The complete codebase developed for this thesis — preprocessing
pipelines, the three interpolation models, evaluation notebooks and
the scripts that build the gridded product — is publicly available in
the project's GitHub repository~\cite{thesisRepo}.
